\chapter{Cel i zakres prac}
\begin{itemize}
    \item Zbadanie fizycznych i logicznych granic wydajności sieci bezprzewodowych.
    \item Dopasowanie protokołów komunikacyjnych do wymagań systemów wbudowanych.
    \item Stworzenie aplikacji analitycznej do oceny jakości transmisji danych w sieciach bezprzewodowych.
    \item Akwizycja danych za pomocą sprzętowej symulacji systemu wbudowanego (HIL).
    \item Interaktywne wizualizacje wyników badań w aplikacji analitycznej.
\end{itemize}



\chapter{Wstęp}

\section{Problematyka systemów czasu rzeczywistego}
\begin{itemize}
    \item sterowanie zdalne
    \item telemetria
    \item determinizm nowoczesnego IoT
\end{itemize}

\section{System analityczny do oceny jakości transmisji danych w sieciach bezprzewodowych}
\begin{itemize}
    \item po co? - dlaczego warto badać protokoły komunikacyjne w systemach wbudowanych?
    \item co? - jakie protokoły komunikacyjne zostaną zbadane?
    \item jak? - w jaki sposób zostaną przeprowadzone badania?
    \item dlaczego? - jakie są oczekiwane rezultaty badań?
    \item dostępne alternatywy
    \item ograniczenia i kompromisy
\end{itemize}



\chapter{Koncepcja Hardware-in-the-Loop (HIL) i autorski system badaczy}
\section{Hardware-in-the-Loop (HIL)}
Czym jest HIL? Dlaczego warto stosować HIL w badaniach systemów wbudowanych? Jakie są zalety i ograniczenia HIL?

Wyjaśnienie zasady działania pętli sprzętowej, w której fizyczne urządzenie "myśli", że steruje realnym obiektem, wymieniając dane z symulatorem.

\section{Architektura}
\subsection{Control Plane}
nawiązywanie połączenia, sterowanie cyklem
\subsection{Data Plane}
brutalne, wysokoczęstotliwościowe tłoczenie pakietów

\section{Akwizycja danych}
Fetch, ekstrakcja i formaty zapisu

\section{Macierz testowa}
\subsection{Częstotliwość transmisji}
\subsection{Topologie sieciowe}
\subsection{Payload}



\chapter{Platforma sprzętowa ESP32}

\section{Architektura mikrokontrolera ESP32}
Opis dwurdzeniowego układu ESP32, jego zasobów sprzętowych, interfejsów radiowych i pamięci operacyjnej.

\section{System czasu rzeczywistego FreeRTOS}
Wyjaśnienie, w jaki sposób wykorzystano system operacyjny czasu rzeczywistego do priorytetyzacji wątków sieciowych i zagwarantowania dokładności pomiarów.

\section{Ekosystem PlatformIO}
Opis nowoczesnego podejścia do programowania systemów wbudowanych, zarządzania zależnościami i kompilacji z użyciem PlatformIO (w kontrze do klasycznego Arduino IDE).



\chapter{Charakterystyka badanych protokołów komunikacyjnych}

\section{Architektura stosu sieciowego (Model OSI i TCP/IP)}
Ogólne wprowadzenie do tego, jak dane są pakowane w kolejne warstwy (od warstwy fizycznej po aplikacyjną) i jaki generuje to narzut (overhead).

\section{Łącze szeregowe (UART)}
Opis najprostszego, niskopoziomowego protokołu, jego zastosowań do twardego sterowania po kablu oraz limitów buforów FIFO.

\section{Protokół UDP (User Datagram Protocol)}
Opis bezpołączeniowej natury protokołu "wyślij i zapomnij" oraz jego kluczowej roli w telemetrii i streamingu wideo/audio.

\section{Protokół TCP (Transmission Control Protocol)}
Wyjaśnienie zasady działania okna przesuwanego i potwierdzeń (ACK). Analiza, dlaczego został stworzony do bezbłędnego transferu plików, a nie do sterowania.

\section{WebSocket (WS)}
Charakterystyka protokołu stałego połączenia nad HTTP. Opis jego zastosowania w nowoczesnych aplikacjach webowych, smart home i interfejsach HMI.

\section{MQTT (Message Queuing Telemetry Transport)}
Przedstawienie standardu IoT opartego na architekturze Pub/Sub z centralnym brokerem. Wyjaśnienie różnych poziomów jakości usług (QoS 0, 1, 2).

\section{Standardy Bluetooth (Classic SPP oraz BLE)}
Porównanie starszego profilu SPP (symulacja kabla) do nowoczesnego, energooszczędnego BLE, z uwzględnieniem pojęcia "Connection Interval" i jego rynkowych zastosowań (od słuchawek po czujniki tętna).



\chapter{Metryki oceny jakości transmisji danych (QoS)}

\section{Skuteczność dostarczania (PDR)}
Matematyczny sposób wyliczania stosunku pakietów odebranych do wysłanych i co ta wartość oznacza w praktyce.

\section{Zmienność opóźnień (Jitter)}
Sposób obliczania fluktuacji czasu między pakietami (odchylenie standardowe, współczynnik CV) jako głównego wskaźnika płynności sterowania.

\section{Zjawisko utraty ciągłej (Burst Loss)}
Definicja i metoda detekcji serii zgubionych pakietów z rzędu (blackoutów sterowania), najczęściej wywoływanych przepełnieniem buforów.

\section{Wydajność użyteczna (Goodput)}
Algorytm odrzucający narzuty sieciowe, wyliczający czystą przepustowość wyłącznie udanych i unikalnych danych z perspektywy odbiorcy.
(Pamiętać o ciekawym przypadku gdzie przekroczono 100 procent)

\section{Defekty kolejności (Out-of-Order - OoO)}
Mechanizm zliczania pakietów, które na skutek specyfiki medium radiowego "wyprzedziły się" w powietrzu i dotarły w błędnej kolejności.

\section{Trendy czasowe i zator buforów (Bufferbloat)}
Obliczanie maksymalnego odchylenia czasu dostarczenia, będącego dowodem na utykanie pakietów w kolejkach rutera lub systemu operacyjnego.

\section{Diagnostyka silnika testowego (Engine Timing)}
Omówienie parametrów logistycznych skryptu: \texttt{engine\_time\_tx\_theory} (czas teoretyczny), \texttt{engine\_time\_tx\_actual} (czas rzeczywistego nadawania), \texttt{engine\_time\_fetch} (czas pobierania) i \texttt{engine\_time\_total\_loop}.

\section{Znormalizowane parametry QoS}
Parametry QoS pozwalające na porównanie różnych protokołów komunikacyjnych w ujednolicony sposób, niezależnie od ich specyfiki i narzutów sieciowych.



\chapter{Architektura aplikacji analitycznej}

\section{Ekosystem Vue.js i framework Nuxt}
Wyjaśnienie wyboru technologii,aplikacja typu SPA (Single Page Application), SSR (Server Side Rendering) i SSG (Static Site Generation), a także zalety i wady tych podejść.

\section{Warstwa wizualna: Tailwind CSS i Nuxt UI}
Krótki opis budowy responsywnego interfejsu użytkownika opartego na nowoczesnych klasach narzędziowych i gotowych komponentach UI.

\section{Wizualizacja danych pomiarowych z Chart.js}
Opis technicznej implementacji interaktywnych wykresów w przeglądarce, zasilanych danymi z wygenerowanej bazy SQLite.



\chapter{Analiza wyników badań i studium przypadków}

\section{Wzorzec odniesienia (Baseline)}
kabel kontra eter - Prezentacja zrzutów ekranu z aplikacji Nuxt ukazujących idealne statystyki dla złącza Serial oraz moment sprzętowego załamania powyżej 1000 Hz.

\section{Brak gwarancji PDR w protokole UDP}
Analiza wykresów udowadniająca, że protokół UDP naturalnie gubi pakiety przy obciążeniu sieci (widoczne spadki PDR poniżej 100\%).

\section{Opóźnienia TCP podczas retransmisji}
Szczegółowe studium przypadku, pokazujące jak protokół TCP za wszelką cenę "dowozi" pakiety, płacąc za to gigantycznym wzrostem Jittera i zjawiskiem Bufferbloatu.

\section{Porównanie topologii i wpływu wielkości Payload}
Wykresy ilustrujące różnicę w zarządzaniu pamięcią przez domowy router w stosunku do ESP32 jako punktu dostępowego oraz analiza momentu wysycenia pasma dla dużych ramek.



\chapter{Wnioski}

\section{Weryfikacja tez badawczych}
Krótkie zebranie dowodów na to, że nie istnieje protokół idealny i inżynier musi zawsze iść na kompromis między opóźnieniem a rzetelnością.

\section{Rekomendacje aplikacyjne}
Praktyczny przewodnik po architekturze na bazie wyników, m.in. zalecenie używania UDP/ESP-NOW do dynamicznego sterowania (samochody RC, drony) oraz TCP/WebSocket do systemów powolnych, bezstratnych (stacje pogodowe, loggery).

\section{Perspektywy dalszego rozwoju projektu}
Stworzenie dedykowanego kontrolera sprzętowego w C++.
